\section{Renew the Heart with St.~Francis de Sales}

St.~Francis de Sales gives the meditation on hell a place within the beginning of the devout life. His sequence leads the soul through the things that disclose its end and the choices which follow from it. In the meditation itself he considers the sufferings of the senses, the greater pain of losing the vision of God's glory, and the eternity of the deprivation. He then asks for definite resolutions by which sin is to be avoided (\emph{Introduction to the Devout Life}, I.15).

Francis moves from the pains of hell to resolutions against sin. An impression of suffering can fade while an attachment remains protected. He asks the person meditating to consider what conduct must change and to resolve upon the means of changing it. This requires honesty about an evil which can seem small when considered only in its immediate attraction. The resolution is strengthened by the thought of what that evil endangers. A person who has considered the everlasting loss of God has reason to accept the inconvenience of removing an occasion of grave sin.

His fifth part turns to the renewal needed after good resolutions have already been made. He compares the heart's devotion to a clock. It requires regular winding and, at intervals, a more careful examination of its workings. Rust must be cleared, what is bent put right, and what is worn attended to. At least once a year, the devout person should examine the affections and passions and renew the purpose of serving God (V.1).

The image recognizes gradual deterioration. A soul may not have made a dramatic decision to abandon prayer. Prayer was shortened, delayed, displaced by other business, and finally omitted. A habit of generosity may have narrowed as comforts multiplied. A fault once resisted may have acquired a convenient explanation. The annual recollection gives time to see developments that are hard to notice in the middle of daily activity.

Francis joins the renewal to Confession and the Eucharist. The examination exposes what needs healing, and the sacraments bring the soul again to the mercy and life of Christ. His counsel is practical about time: choose a suitable occasion, seek greater recollection, and make the meditations with the advice of one's spiritual father. The end of the liturgical year can provide such an occasion, especially when the remembrance of death has already made the soul attentive.

He begins the renewal with the benefits of God's call. The person remembers the renunciation of mortal sin, the dedication of the whole self to divine service, and the resolve to rise again through grace if a fall should occur (V.2). He considers the helps by which God drew him: prayer, sacraments, spiritual reading, and the growth of love. Conversion is remembered as mercy, and progress is received with thanksgiving.

Francis also attends to those called late in life. The years lost do not make the present call worthless. Mercy has reached the sinner before death and has opened a way of service. That remembrance places regret within gratitude. We can grieve over time wasted while giving thanks that the time now offered can be used differently.

The examination that follows is to be made quietly and patiently. Francis permits its parts to be distributed across an appropriate period; he does not require the whole to be performed in one strained sitting. He directs attention toward conduct in relation to God, oneself, one's neighbor, and the passions. Progress is to produce thanksgiving without self-congratulation. Little progress, or even decline, is to produce renewed humility and effort with God's help, rather than discouragement (V.3).

At the close, Francis names the means that must support the renewed purpose: prayer, sacraments, good works, correction of faults already discovered, avoidance of occasions, and fidelity to counsel (V.15). The annual examination can show where one of those means has been allowed to lapse. If prayer has repeatedly been displaced, merely resolving to pray more leaves the old difficulty untouched. A suitable time must be set aside and kept when competing demands return. The resolution then addresses the way devotion actually weakened, instead of depending upon the recollection's strong impression to continue.

Francis's comparison with a clock concerns continuing attention to devotion. A yearly examination gives time to discover a fault which ordinary haste has concealed; the means he names help the person correct it afterward. The resolution must fit the duty and the weakness discovered. An indiscriminate increase of practices can leave the neglected obligation as neglected as before. The person should use counsel to choose a definite amendment he can faithfully undertake, remembering that the time for serving God remains limited. Renewal asks him to begin that amendment within his present vocation.

